\documentclass[11pt]{article}

\usepackage[margin=1in]{geometry}
\usepackage{amsmath}
\usepackage{amssymb}
\usepackage{booktabs}
\usepackage{array}
\usepackage{enumitem}
\usepackage{xcolor}
\usepackage{hyperref}
\hypersetup{
  colorlinks=true,
  linkcolor=black,
  urlcolor=blue,
  citecolor=black
}

\newcommand{\tvec}{\mathbf{t}}
\newcommand{\ovec}{\mathbf{o}}
\newcommand{\norm}[1]{\left\lVert #1 \right\rVert}

\title{An Envelope-Fit Metric for Charge-State and Monoisotope Selection}
\author{FlashLFQ-Rust}
\date{\today}

\begin{document}
\maketitle

\begin{abstract}
In untargeted MS1 feature detection, deciding which \emph{charge state} and
monoisotopic mass best explain an observed isotopic envelope is a model-selection
problem. The current detector scores a hypothesis with an un-normalized weighted
sum of comb-weighted, retention-time--smoothed intensities. That statistic rewards
matched peaks but is blind to two failure modes: predicted isotope teeth that are
\emph{missing}, and observed peaks the hypothesis leaves \emph{unexplained}. As a
consequence it systematically mislabels a real charge-$2$ envelope (isotopes
$0.5$~Th apart) as charge-$1$ (isotopes $1.0$~Th apart), because the $z{=}1$ comb
lands on every other peak and the raw sum does not care about the peaks in between.
We propose a single unified score: the \emph{cosine similarity} between the
theoretical averagine envelope and the observed intensities, evaluated over the
\emph{union} of predicted-isotope positions and in-window observed peaks. We show
that this one quantity simultaneously measures goodness of fit, the fraction of
observed signal explained, and predicted-peak completeness --- and that its square
is exactly the fraction of observed signal energy captured by the best-scaled
envelope. On a concrete peptide (EFLDVIK, true charge $2$) the metric cleanly ranks
the true charge above both the harmonic $z{=}1$ and the over-split $z{=}3$
hypotheses, where no single sub-metric does.
\end{abstract}

\section{Motivation}
\label{sec:motivation}

\subsection{The \texttt{RawSum} statistic and why it fails}

For a charge/monoisotope hypothesis the detector currently accumulates
\begin{equation}
\texttt{RawSum} \;=\; \sum_{k}\sum_{s} w_k \, g_s \, I_{k,s},
\label{eq:rawsum}
\end{equation}
where $k$ indexes the predicted isotope teeth of the hypothesis, $s$ indexes scans
across the retention-time (RT) apex, $w_k$ is the averagine comb weight of tooth
$k$, $g_s$ is an RT-Gaussian weight centered on the apex, and $I_{k,s}$ is the
observed intensity found at the $m/z$ of tooth $k$ in scan $s$.

Equation~\eqref{eq:rawsum} is a \emph{matched-filter response}. It grows with every
predicted tooth that happens to find intensity, and it is never decremented. Two
things it cannot see:
\begin{enumerate}[label=(\alph*)]
  \item \textbf{Missing predicted teeth.} If a tooth the model requires has no
  intensity, that tooth simply contributes $0$ to the sum. A hypothesis that
  explains half of its own predicted envelope is not penalized relative to one that
  explains all of it, as long as the absolute matched intensity is comparable.
  \item \textbf{Unexplained observed peaks.} Any observed peak that does \emph{not}
  sit on a predicted tooth is invisible to Eq.~\eqref{eq:rawsum} entirely. A
  hypothesis is free to ignore half the signal in the window at no cost.
\end{enumerate}
Because the score is neither normalized for completeness nor penalized for the
signal it ignores, it does not distinguish a correct charge from a \emph{harmonic}
charge.

\subsection{The harmonic example}

Consider a genuine charge-$2$ peptide. Its isotope teeth are spaced
$1.00\text{--}\mathrm{Da}/2 \approx 0.5$~Th apart. A charge-$1$ hypothesis predicts
teeth $1.0$~Th apart --- i.e. on \emph{every other} observed peak. Schematically,
over eight observed peaks $P_0,\dots,P_7$ spaced $0.5$~Th:

\begin{center}
\ttfamily
\begin{tabular}{l l}
observed ($0.5$~Th grid): & P0  P1  P2  P3  P4  P5  P6  P7 \\
$z{=}2$ comb (true):      & \^{}\ \ \^{}\ \ \^{}\ \ \^{}\ \ \^{}\ \ \^{}\ \ \^{}\ \ \^{}\ \ (hits every peak) \\
$z{=}1$ comb (harmonic):  & \^{}\ \ \ \ \ \ \^{}\ \ \ \ \ \ \^{}\ \ \ \ \ \ \^{}\ \ \ \ (hits every \emph{other} peak) \\
\end{tabular}
\end{center}

The $z{=}1$ comb finds real intensity at $P_0, P_2, P_4, P_6$ and produces a large
\texttt{RawSum}. The intervening peaks $P_1, P_3, P_5, P_7$ --- which are precisely
the evidence that the true spacing is $0.5$~Th, not $1.0$~Th --- contribute nothing
to the $z{=}1$ score and nothing against it. The matched filter is happy; the charge
label is wrong. What is missing is any term that charges the $z{=}1$ hypothesis for
the half of the envelope it declines to explain.

\section{The three requirements}
\label{sec:requirements}

A score fit for charge/monoisotope selection must combine three properties that the
harmonic example shows to be individually insufficient:
\begin{enumerate}[label=\textbf{P\arabic*.},leftmargin=3.2em]
  \item \textbf{Goodness of fit.} The observed intensities should match the
  \emph{shape} of the theoretical isotope-ratio pattern, not merely be present.
  \item \textbf{Explained fraction.} Observed signal inside the window that the
  hypothesis does not predict should count against it; a hypothesis may not ignore
  half the peaks for free.
  \item \textbf{Completeness.} A predicted tooth that is absent from the data should
  count against the hypothesis; a half-present envelope is a worse explanation than
  a fully present one.
\end{enumerate}
The central claim of this note is that these three requirements are not three terms
to be weighted and summed. They are three behaviors of one quantity.

\section{The metric}
\label{sec:metric}

\subsection{Union grid and formula}

Let the \emph{window} be a fixed $m/z$ interval around the candidate monoisotope
(Section~\ref{sec:design}). Define two index sets:
\begin{align*}
\mathcal{T} &= \{\text{predicted isotope positions of the hypothesis with weight above the significance cutoff}\},\\
\mathcal{O} &= \{\text{observed peaks in the window above the noise floor}\}.
\end{align*}
Form the \textbf{union grid} $\mathcal{G} = \mathcal{T} \cup \mathcal{O}$ by merging
positions that coincide (within an $m/z$ tolerance) and keeping the rest. On this
common grid build two vectors indexed by $j \in \mathcal{G}$:
\[
t_j = \begin{cases} \text{theoretical averagine weight at } j, & j \in \mathcal{T},\\[2pt] 0, & j \notin \mathcal{T},\end{cases}
\qquad
o_j = \begin{cases} \text{observed intensity at } j, & j \in \mathcal{O},\\[2pt] 0, & j \notin \mathcal{O}.\end{cases}
\]
The score is the cosine similarity of these two vectors:
\begin{equation}
\boxed{\;
S \;=\; \frac{\sum_{j\in\mathcal{G}} t_j\, o_j}
             {\sqrt{\sum_{j\in\mathcal{G}} t_j^2}\;\sqrt{\sum_{j\in\mathcal{G}} o_j^2}}
   \;=\; \frac{\tvec\cdot\ovec}{\norm{\tvec}\,\norm{\ovec}}\; }
\label{eq:cosine}
\end{equation}
with $S \in [0,1]$ for non-negative intensities.

\subsection{How one formula delivers all three properties}

The union grid is what makes Eq.~\eqref{eq:cosine} do three jobs at once. Every grid
slot falls into one of three cases, and each case moves $S$ in exactly the direction
one of P1--P3 demands.

\begin{table}[h]
\centering
\renewcommand{\arraystretch}{1.35}
\begin{tabular}{@{}p{3.6cm} c c p{5.4cm}@{}}
\toprule
\textbf{Grid slot} & $t_j$ & $o_j$ & \textbf{Effect on $S$} \\
\midrule
Matched tooth
  & large & large
  & Raises the numerator $\tvec\cdot\ovec$ when the ratio $o_j/t_j$ agrees with the
    envelope shape \textbf{(P1: fit)}. \\
Missing predicted tooth ($j\in\mathcal{T}\setminus\mathcal{O}$)
  & large & $0$
  & Contributes to $\norm{\tvec}$ (denominator) but nothing to the numerator, so $S$
    drops \textbf{(P3: completeness)}. \\
Unexplained observed peak ($j\in\mathcal{O}\setminus\mathcal{T}$)
  & $0$ & large
  & Contributes to $\norm{\ovec}$ (denominator) but nothing to the numerator, so $S$
    drops \textbf{(P2: explained fraction)}. \\
\bottomrule
\end{tabular}
\caption{Each type of union-grid slot maps to one required property. There are no
free parameters trading the three off; they are the same denominator/numerator
bookkeeping.}
\label{tab:mapping}
\end{table}

Concretely, in the harmonic example the intervening peaks $P_1,P_3,P_5,P_7$ enter
the $z{=}1$ score as $(t_j = 0,\, o_j = I)$ slots. They inflate $\norm{\ovec}$ with
no numerator contribution and pull the $z{=}1$ cosine down --- exactly the penalty
\texttt{RawSum} was missing.

\subsection{Why $S^2$ \emph{is} the explained-signal fraction}

The ``portion of signal explained'' is not a bolted-on term; it is what
Eq.~\eqref{eq:cosine} already computes. Model the observed vector as a scaled copy of
the envelope plus residual, $\ovec \approx A\,\tvec$, and choose the amplitude $A$ by
least squares:
\begin{equation}
A^\star = \arg\min_A \norm{\ovec - A\tvec}^2
        = \frac{\tvec\cdot\ovec}{\norm{\tvec}^2}.
\label{eq:amp}
\end{equation}
The energy of the observed signal splits into the part captured by the best-scaled
envelope and the residual. Using $A^\star$,
\begin{align}
\norm{\ovec - A^\star \tvec}^2
  &= \norm{\ovec}^2 - 2A^\star (\tvec\cdot\ovec) + (A^\star)^2\norm{\tvec}^2 \nonumber\\
  &= \norm{\ovec}^2 - \frac{(\tvec\cdot\ovec)^2}{\norm{\tvec}^2}.
\label{eq:residual}
\end{align}
Define the coefficient of determination through the origin,
$R^2 = 1 - \dfrac{\norm{\ovec - A^\star\tvec}^2}{\norm{\ovec}^2}$
(the fraction of observed energy the scaled envelope removes). Substituting
Eq.~\eqref{eq:residual},
\begin{equation}
R^2 = \frac{(\tvec\cdot\ovec)^2}{\norm{\tvec}^2\,\norm{\ovec}^2} = S^2 .
\label{eq:r2}
\end{equation}
So $S^2$ is \emph{literally} the fraction of observed signal energy explained by the
best-scaled averagine envelope, evaluated over a grid that already contains the
unexplained peaks. Property P2 is therefore not an add-on: once the union grid
includes the observed peaks the envelope ignores, the cosine's square is the
explained-energy fraction. ($S$ and $S^2$ are monotonically related, so the choice
between reporting the cosine or its square is one of interpretation only.)

\section{Design choices}
\label{sec:design}

Four judgment calls determine the metric's behavior. They are where the real
engineering lives.

\paragraph{The window (what ``should be explained'' means).}
The window sets the domain of $\mathcal{O}$ --- the set of observed peaks the
hypothesis is held responsible for. We take it to span from the candidate
monoisotope to $\text{mono} + n\cdot\Delta$ (where $\Delta$ is the isotope spacing
and $n$ the number of significant teeth) plus a small margin. \textbf{The same
window, in $m/z$, is used for every charge hypothesis at a given monoisotope.} This
is deliberate: a co-eluting interferent inside the window then penalizes every charge
hypothesis equally, so it drags all absolute scores down together and leaves the
\emph{comparison} between charges fair. Using a per-charge window would let a
hypothesis quietly shrink its own accountability. The window boundary is the most
consequential parameter in the whole scheme.

\paragraph{Noise threshold.}
Before an observed peak is added as a $(t=0,\,o=I)$ slot, it must clear a noise
floor $\eta$. Without this, ordinary baseline noise in the window would read as
``unexplained signal'' and depress every cosine, muddying the comparison and
biasing the metric against the true (fully-explaining) hypothesis. Peaks below
$\eta$ are simply not placed on the grid.

\paragraph{Cosine, not Pearson.}
Use cosine similarity, \emph{not} the Pearson correlation. Pearson mean-centers each
vector before the dot product; that subtraction destroys the meaning of a zero. An
\emph{absent} isotope peak should read as the number $0$ (``no intensity''), but
Pearson re-interprets it as ``below the average intensity,'' which is not the same
thing and can even reward a hypothesis for a systematically empty tooth. The zero
must mean zero for P2 and P3 to work, so mean-centering is disqualifying here.
(FlashLFQ uses Pearson correlation for its binary accept/reject \emph{gate}; that is
a different task. For a fit-and-coverage \emph{score} over a union grid with
meaningful zeros, cosine is the correct choice.)

\paragraph{Significance cutoff.}
Which theoretical teeth count as ``predicted'' (the set $\mathcal{T}$) is set by a
weight cutoff --- e.g.\ keep teeth whose averagine weight is $\ge 0.2$ of the modal
tooth. Teeth below the cutoff are negligible in the theoretical pattern; including
them would add near-zero $t_j$ slots that only add noise to the completeness
accounting.

\section{Validation}
\label{sec:validation}

We evaluated the metric on peptide \textbf{EFLDVIK}, true charge $2$, monoisotopic
mass $862.488$~Da, with $15$ observed peaks inside the window. The existing detector
had \emph{mislabeled} this feature as charge $1$ at mass $431$~Da --- the harmonic
error of Section~\ref{sec:motivation}. Table~\ref{tab:validation} reports, for each
charge hypothesis, completeness (fraction of predicted teeth present), the explained
fraction (\%), and the union-grid cosine $S$.

\begin{table}[h]
\centering
\renewcommand{\arraystretch}{1.25}
\begin{tabular}{@{}l c c c c@{}}
\toprule
\textbf{Hypothesis} & \textbf{Mass (Da)} & \textbf{Completeness} & \textbf{\% explained} & \textbf{cosine $S$ (union)} \\
\midrule
$z=1$ (harmonic)      & $431$  & $1.00$ & $0.60$ & $0.905$ \\
$\mathbf{z=2}$ (true) & $862$  & $1.00$ & $0.76$ & $\mathbf{0.984}$ \\
$z=3$ (over-split)    & $1294$ & $0.33$ & $0.54$ & $0.718$ \\
\bottomrule
\end{tabular}
\caption{EFLDVIK charge selection. The unified cosine ranks the true charge first;
neither sub-metric does so alone.}
\label{tab:validation}
\end{table}

\paragraph{What the table shows.}
Read column by column, the sub-metrics each fail on their own:
\begin{itemize}[leftmargin=1.4em]
  \item \textbf{Completeness alone ties $z{=}1$ and $z{=}2$} ($1.00$ vs $1.00$). Both
  the true charge and its harmonic have all their predicted teeth present --- of
  course they do, since the $z{=}1$ teeth are a subset of the $z{=}2$ peaks.
  Completeness cannot break the harmonic tie and cannot fix the mislabel. It does,
  correctly, sink $z{=}3$ ($0.33$), which predicts teeth at $0.33$~Th spacings that
  are largely absent.
  \item \textbf{Explained fraction alone catches $z{=}1$} ($0.60 < 0.76$), because
  the harmonic ignores the intervening peaks --- but on its own it does not cleanly
  separate the field ($z{=}3$ at $0.54$ is close to $z{=}1$ at $0.60$, for the
  unrelated reason that $z{=}3$ predicts too much).
  \item \textbf{The unified cosine separates all three.} $z{=}1$ is pulled down to
  $0.905$ by the unexplained intervening peaks (P2); $z{=}3$ is pulled down to
  $0.718$ by its missing teeth (P3); $z{=}2$ retains a high $0.984$ because it both
  explains the observed peaks and has all its predicted teeth present (P1--P3
  together). The true charge wins.
\end{itemize}
This is the argument for a single cosine-over-union score rather than any one
sub-metric or an ad~hoc weighted blend: the three requirements are already
co-resident in Eq.~\eqref{eq:cosine}, and only their combination ranks the
hypotheses correctly.

\section{Caveats and implementation note}
\label{sec:caveats}

\paragraph{Caveats (stated honestly).}
\begin{itemize}[leftmargin=1.4em]
  \item \textbf{The winning margin is modest.} $0.984$ vs $0.905$ is a real
  separation but not a chasm. Reliable charge \emph{re-selection} across a full
  dataset needs empirical validation before it replaces the existing decision: on
  genuinely noisy or sparse envelopes the metric could begin to prefer higher
  charges (which predict more, finer teeth and can accrue spurious matched
  intensity). This must be measured, not assumed.
  \item \textbf{The window boundary dominates.} As noted in
  Section~\ref{sec:design}, the choice of window is the single most consequential
  parameter; it defines what counts as ``unexplained'' and therefore the entire P2
  penalty. It deserves the most tuning and the most scrutiny.
  \item \textbf{Chimeric envelopes lower the cosine --- correctly.} When two
  features co-elute and overlap, no single-envelope hypothesis explains all the
  in-window signal, and every cosine is depressed. This is the \emph{right} signal:
  a uniformly low $S$ across all charges is a flag to route the window to advanced
  multi-envelope deconvolution rather than to force a single-charge label.
\end{itemize}

\paragraph{Implementation note.}
The detector already contains most of this. The opt-in scorer
\texttt{ScoreModel::NormalizedNoiseFloor} with the \texttt{score\_cosine} flag
computes a cosine, but currently only over the \emph{comb} slots --- i.e.\ over
$\mathcal{T}$ alone. Over $\mathcal{T}$ the cosine already delivers P1 (fit) and P3
(completeness, via the missing-tooth $(t{>}0,o{=}0)$ slots), but it is blind to P2
because unexplained observed peaks off the comb are never on the grid. The single
change that upgrades it to the full three-property metric is to add the in-window
observed peaks that are \emph{not} on the comb as $(t=0,\,o=I)$ slots (after the
noise-floor filter) before taking the cosine. That one addition to the grid turns the
existing comb cosine into Eq.~\eqref{eq:cosine}, and with it $S^2$ becomes the
explained-energy fraction of Section~\ref{sec:metric}.

\end{document}
